\documentclass[12pt]{article}

\usepackage[margin=1in]{geometry}
\usepackage[T1]{fontenc}
\usepackage{parskip}
\usepackage{hyperref}
\hypersetup{
    colorlinks=true,
    linkcolor=black,
    urlcolor=cyan,
    pdftitle={Cover Letter -- Harvard University},
}
\urlstyle{same}

\pagenumbering{gobble}

\begin{document}

\begin{center}
  {\Large\textsc{Edward A. Varvak}}\\[10pt]
\end{center}

Department of Mathematics \\
University of Chicago \\
5734 S University Ave \\
Chicago, IL 60637

\today

Dear Members of the Hiring Committee,

I am a graduate student at the University of Chicago, working in complex
algebraic geometry and Hodge theory under the direction of Benjamin Bakker,
and I expect to receive my Ph.D.\ in May 2027. I am applying for the position
of Benjamin Peirce Fellow in Mathematics (\mbox{Harvard-BPFELLOW20272028}).
My research concerns the monodromy and Hodge theory of Lagrangian fibrations
on hyper-K\"ahler and other symplectic varieties, as well as moduli spaces of
abelian varieties and their Mumford--Tate groups. At Harvard, I would be
particularly interested in working with Mihnea Popa and Mark Kisin.

My main goal in teaching is for students to leave a course able to explain
why the methods they learned work, not only how to carry them out. As such, I try
to show them what a formula is doing before they use it. For example, when 
teaching calculus at the University of Chicago, I go through volumes of revolution
by drawing the solid, sketching a typical slice, and working out its
area with the class, instead of handing out the washer and shell formulas.

When I work through a proof or an example, I make a point of
aking the class what the next step should be and waiting for a guess before going on. 
I find that doing so creates an environment for the class where the students are 
unafraid to ask questions of their own accord.
I also try to match the material I cover to the level of the class. 
At the University of Illinois Chicago, most of my students in
Finite Mathematics for Business took it to satisfy a requirement, and many
arrived unsure of their mathematical ability; with them, I set
aside more advanced applications in favor of making sure every student could
follow.

Outside the classroom, my advising so far has been informal. In the learning seminars 
I organized, I helped speakers prepare their talks. While I would be happy to continue 
organizing seminars at Harvard, I would also like to try out a
more concrete mentoring arrangement, such as advising undergraduate reading projects 
in algebraic geometry. Likewise, I would welcome the chance to teach a one-term 
course in algebraic geometry or Hodge theory.

Along with this letter, I include my curriculum vitae, research statement,
and service statement. Research letters will come from Benjamin Bakker
(Chicago), Radu Laza (Stony Brook), \.{I}zzet Co\c{s}kun (UIC), and Valentino
Tosatti (Courant), and a teaching letter from McKinley Meyer (UIC). Thank you
for your consideration.

Sincerely,

Edward A. Varvak

\end{document}
